%% kapa: exam=1 round=36 year=2025 period=2 form=A subject=appraisal_law subject_ko=감정평가관계법규 seq=9 number=9
%% source: Q-Net 감정평가사 기출문제 (cst00309) artlSeq=5231525 file=raw/1008/5231525/2025년 제36회 감정평가사 1차시험 2교시 문제지 원본.pdf
%% conversion: kordoc 4.24.51 -> tools/gen1.py (auto); review=proofread-sonnet; flags=none
\begin{question}{9}
국토의 계획 및 이용에 관한 법령상 성장관리계획에 관한 설명으로 옳은 것은?
\begin{choices}
  \item 공업지역 중 주변지역과 연계하여 체계적인 관리가 필요한 지역은 성장관리계획구역으로 지정할 수 있다.
  \item 성장관리계획구역 내 보전녹지지역에서는 성장관리계획으로 정하는 바에 따라 50퍼센트 이하로 건폐율을 완화하여 적용할 수 있다.
  \item 성장관리계획구역에서 개발행위 또는 건축물의 용도변경을 하려면 그 성장관리계획에 맞게 하여야 한다.
  \item 성장관리계획구역의 면적을 5퍼센트 이내에서 변경하려면 미리 주민과 지방의회의 의견을 들어야 한다.
  \item 군수는 성장관리계획구역의 지정 또는 변경에 관한 공고를 한 때에는 성장관리계획구역안을 10일 이내로 일반이 열람할 수 있도록 해야 한다.
\end{choices}
\end{question}
\answer{3}
